\documentclass[11pt]{amsart}
\usepackage{amsmath,amssymb,amsthm}
\usepackage{hyperref}
\usepackage{booktabs}
\usepackage[margin=1in]{geometry}

\theoremstyle{plain}
\newtheorem{theorem}{Theorem}[section]
\newtheorem{lemma}[theorem]{Lemma}
\newtheorem{corollary}[theorem]{Corollary}
\newtheorem{proposition}[theorem]{Proposition}
\theoremstyle{definition}
\newtheorem{definition}[theorem]{Definition}
\newtheorem{remark}[theorem]{Remark}
\newtheorem{example}[theorem]{Example}
\theoremstyle{remark}
\newtheorem*{conjecture}{Conjecture}
\newtheorem*{question}{Question}
\newtheorem*{openproblem}{Open Problem}

\newcommand{\Z}{\mathbb{Z}}
\newcommand{\Q}{\mathbb{Q}}
\newcommand{\C}{\mathbb{C}}
\newcommand{\R}{\mathbb{R}}
\newcommand{\pprod}{p_{\mathrm{prod}}}

\title[Modular equations of degrees 11, 13, 17, and a conjecture for general prime degree]{The Degree-Eleven, Degree-Thirteen, and Degree-Seventeen Modular Equations,\\ and a Conjecture for General Prime Degree,\\ in Ramanujan's Prime-Degree Theta Function Theory}
\author{Ashwin Sathya Raghunathan}
\date{\today}
\subjclass[2020]{11F27, 11F11, 33D10, 11Y40}
\keywords{Ramanujan theta functions, modular equations, eta-quotients, Sturm bound, Atkin--Lehner involution}

\begin{document}

\begin{abstract}
Ramanujan recorded, in a marginal note at the end of Section 12 of Chapter 20 of
his second notebook, that the theta functions $\varphi(q^{1/n})$, $\psi(q^{1/n})$,
$f(-q^{1/n})$ for prime $n$ can be expressed as a sum of $\frac12(n-1)$ $n$th roots
plus $\varphi(q^n)$, and that the auxiliary equations needed are quadratics only for
degrees $5, 17, 257$, cubics for $7, 13, 19, 37, 73, 97, 109, 163, 193$, and quintics
for $11, 41, 101$. The cubic, quintic, and septic cases of this envisioned theory
have been proved in recent work of Berndt--Reb\'ak and Reb\'ak; no case for any prime
$p \geq 11$ was known. We find and prove the degree-$11$ modular equation: with
$m = \varphi^2(q)/\varphi^2(q^{11})$ and $\pprod = 32q^5\chi(q)/\chi^{11}(q^{11})$,
writing $Y = m^2$ and $P = \pprod/32$, the explicit relation $G(Y,P) = 0$ holds,
a plane quintic of genus $2$ with small integer coefficients whose pure parts are
the perfect fifth powers $(1-Y)^5$ and $(1+32P)^5$; it is proved by the Sturm-bound
method on $\Gamma_0(44)$. We prove the next two cases by the same method: the
degree-$13$ modular equation $F(m,P) = 0$, a plane septic of genus $2$ proved on
$\Gamma_0(52)$, and the degree-$17$ modular equation $F(m,P) = 0$, a plane curve of
degree $12$ and genus $3$ proved on $\Gamma_0(68)$. In all three cases the pure parts
of the relation are perfect powers with exponents equal to the degrees, and the map
$X_0(4p) \to C$ has degree $2$. We also prove a general product formula for the product
of the dissection terms, valid for every odd $n$, by an elementary argument from the
Jacobi triple product; verify the dissection, the three flagship evaluations, and
eleven class invariants computationally to 50 digits; give Sturm-bound proofs of the
multiplier relations of degrees $3$, $5$, and $7$; and compute the class invariant
$G_{121}$ numerically to 60 digits. The six proved cases suggest a uniform
conjectural template for general prime $p$, driven by the Atkin--Lehner involution
$w_4$ of $X_0(4p)$, with testable predictions for $p = 19$ and $p = 23$.
The auxiliary quintic $r_{11}$ remains open, and
we give a concrete attack plan for it.
\end{abstract}

\maketitle

\section{Introduction}
\label{sec:intro}

At the end of Section 12 of Chapter 20 of his second notebook, Ramanujan wrote a
marginal note (Berndt, \emph{Ramanujan's Notebooks, Part III}, p.~400) asserting that,
for prime $n$, the theta functions $\varphi(q^{1/n})$, $\psi(q^{1/n})$ and
$f(-q^{1/n})$ can be expressed as a sum of $\frac12(n-1)$ $n$th roots of several
functions plus $\varphi(q^n)$, $\psi(q^n)$ or $f(-q^n)$, and that ``in finding the
values of these functions, quadratics only appear in the case of the 5th, 17th,
257th, etc.\ degrees, and cubics in case of the 7th, 13th, 19th, 37th, 73rd, 97th,
109th, 163rd, 193rd etc.\ degrees \dots\ and quintics in case of the 11th, 41st,
101st etc.\ degrees.'' Berndt writes: ``We are unable to provide a proper
interpretation for most of this statement.''

The three cases Ramanujan's note points to have now been proved in full:
\begin{itemize}
\item Degree $3$ (cubic): Berndt--Reb\'ak \cite{BR25}, Theorem 2.1.
\item Degree $5$ (quintic): Berndt--Reb\'ak \cite{BR25}, Theorem 4.2.
\item Degree $7$ (septic): Ramanujan's Lost Notebook, p.~206, Entry 1.1(i)--(v)
proved by Son \cite{Son01}; part (vi), the three missing terms, completed by
Reb\'ak \cite{Reb23}, Theorem 4.1.
\item Degree $9$ (nonic, composite): Kim--Reb\'ak \cite{KR25}.
\end{itemize}
Berndt--Reb\'ak call these ``special cases of a grand theory that Ramanujan
envisioned'' in the note, and warn that ``as the order increases, we know less and
less about the modular equations necessary for a possible grand theory.'' No case
for any prime $p \geq 11$ had been published before this work; we now prove the
cases $p = 11$, $13$, and $17$.

The prime lists in Ramanujan's note admit a clean arithmetic reading, which we
record here since it shapes everything that follows: the lists are exactly the
primes for which $(p-1)/2$ has prime factors only in $\{2\}$, $\{2,3\}$, $\{2,5\}$
respectively. That is, they are the primes whose real cyclotomic field
$\Q(\zeta_p)^+$, of degree $(p-1)/2$ with cyclic Galois group, can be built by
successively adjoining roots of equations of degree $2$ only, $2$ and $3$, or $2$
and $5$. The lists are complete for the stated bound: for example $p=23$ has
$(p-1)/2 = 11$, needing 11th roots, and $p=29$ has $(p-1)/2 = 14 = 2\cdot 7$,
needing 7th roots, so neither appears in any list.

This paper has five contributions.
\begin{enumerate}
\item \textbf{The degree-$11$ modular equation} (Theorem~\ref{thm:p11}): the explicit
relation $G(Y,P) = 0$ between $Y = \varphi^4(q)/\varphi^4(q^{11})$ and
$P = \pprod/32$, proved by the Sturm-bound method on $\Gamma_0(44)$. This is the
first case of Ramanujan's envisioned theory for a prime $p \geq 11$. The relation is
a plane quintic of genus $2$ with small integer coefficients, whose pure parts are
the perfect fifth powers $(1-Y)^5$ and $(1+32P)^5$.
\item \textbf{The degree-$13$ and degree-$17$ modular equations}
(Theorems~\ref{thm:p13} and~\ref{thm:p17}, \S\ref{sec:p1317}): the explicit relations
$F(m,P) = 0$ for $p = 13$ (a plane septic of genus $2$, proved on $\Gamma_0(52)$)
and $p = 17$ (a plane curve of degree $12$ and genus $3$, proved on $\Gamma_0(68)$),
found and proved by the same method as Theorem~\ref{thm:p11}. In both cases the
pure parts are perfect powers with exponents equal to the degrees, and the map
$X_0(4p) \to C$ has degree $2$.
\item \textbf{A general product formula} (Theorem~\ref{thm:pprod}): the product of
the dissection terms is an explicit $\chi$-quotient for \emph{every} odd $n$,
proved by an elementary argument from the Jacobi triple product. This unifies the
three known cases and gives the $p=11$ product used in (1).
\item \textbf{Computational verification} (\S\ref{sec:comp}): the dissection, the
three flagship evaluations, and eleven class invariants verified to 50 digits; the
class invariant $G_{121}$ computed numerically to 60 digits (no closed form is
known).
\item \textbf{Sturm-bound proofs} (\S\ref{sec:sturm}): the multiplier relations of
degrees $3$, $5$, and $7$ proved by the modern route (Newman--Ligozat criteria plus
a Sturm-bound coefficient check), with a new prover; and the $p=11$ frontier
(\S\ref{sec:p11}): the auxiliary quintic $r_{11}$ remains open, and we give a
concrete attack plan for it.
\end{enumerate}

\subsection{Notation}
We use Ramanujan's notation throughout. For $|ab|<1$,
\begin{equation}
f(a,b) := \sum_{n=-\infty}^{\infty} a^{n(n+1)/2} b^{n(n-1)/2}
= (-a;ab)_\infty(-b;ab)_\infty(ab;ab)_\infty,
\end{equation}
the Jacobi triple product. The three principal specialisations are
\begin{equation}
\varphi(q) := f(q,q) = \sum_{n=-\infty}^{\infty} q^{n^2}, \qquad
\psi(q) := f(q,q^3) = \sum_{n=0}^{\infty} q^{n(n+1)/2}, \qquad
\chi(q) := (-q;q^2)_\infty.
\end{equation}
For $q = e^{-\pi\sqrt{n}}$, $n>0$ rational, the class invariant is
\begin{equation}
G_n := 2^{-1/4} q^{-1/24}\chi(q), \qquad G_n = G_{1/n}.
\end{equation}
The multiplier of degree $d$ is $m := \varphi^2(q)/\varphi^2(q^d)$ and we write
$Y := \varphi^4(q)/\varphi^4(q^p) = m^2$. Throughout, $0<q<1$ is real unless
otherwise stated, and $n$ denotes an odd positive integer.

\section{The dissection and the general product formula}
\label{sec:pprod}

The universal dissection (Entry 31, Chapter 16; Kim--Reb\'ak \cite{KR25}, Lemma 2.3)
holds for every odd $n$:
\begin{equation}
\frac{\varphi(q^{1/n})}{\varphi(q^n)}
= 1 + \sum_{k=1}^{(n-1)/2} \frac{2q^{k^2/n}\, f(q^{n+2k},\, q^{n-2k})}{\varphi(q^n)}.
\label{eq:dissection}
\end{equation}
For prime $p$ we define the dissection terms
\begin{equation}
u_k := \frac{2q^{k^2/p}\, f(q^{p+2k},\, q^{p-2k})}{\varphi(q^p)},
\qquad k = 1, \dots, \tfrac{p-1}{2},
\label{eq:uk}
\end{equation}
so that $\varphi(q^{1/p})/\varphi(q^p) = 1 + \sum_k u_k$. For $0<q<1$ the $u_k$
are positive and strictly decreasing in $k$ (Reb\'ak \cite{Reb23}, Lemma 3.2).
We write $\pprod := \prod_k u_k$ for their product.

The three proved cases give:
\begin{center}
\begin{tabular}{c|c|c}
$p$ & $\pprod$ & in class invariants \\ \hline
3 & $2q^{1/3}\chi(q)/\chi^3(q^3)$ & $\sqrt{2}\,G_n/G_{9n}^3$ \\
5 & $4q\,\chi(q)/\chi^5(q^5)$ & $2\,G_n/G_{25n}^5$ \\
7 & $8q^2\chi(q)/\chi^7(q^7)$ & $2\sqrt{2}\,G_n/G_{49n}^7$
\end{tabular}
\end{center}
These fit a single formula, which we now prove for every odd $n$.

\begin{theorem}[General product formula]
\label{thm:pprod}
Let $n$ be odd and $|q|<1$. With $u_k$ as in \eqref{eq:uk} (with $p$ replaced by
$n$), the product of the dissection terms is
\begin{equation}
\prod_{k=1}^{(n-1)/2} u_k
= 2^{(n-1)/2}\, q^{(n^2-1)/24}\, \frac{\chi(q)}{\chi^n(q^n)}.
\label{eq:pprod-chi}
\end{equation}
In particular, for $q = e^{-\pi\sqrt{n}}$,
\begin{equation}
\pprod = 2^{(n-1)/4}\, \frac{G_n}{G_{n^2}^{\,n}}.
\label{eq:pprod-G}
\end{equation}
\end{theorem}

\begin{proof}
By the Jacobi triple product,
\begin{equation}
f(q^{n+2k}, q^{n-2k})
= (-q^{n+2k};q^{2n})_\infty\,(-q^{n-2k};q^{2n})_\infty\,(q^{2n};q^{2n})_\infty.
\end{equation}
As $k$ runs over $1, \dots, (n-1)/2$, the exponents $n+2k$ run over the odd residues
$n+2, n+4, \dots, 2n-1$ modulo $2n$, and $n-2k$ run over the odd residues
$1, 3, \dots, n-2$ modulo $2n$. Together they run over all odd residues modulo $2n$
except $n$ itself. Hence
\begin{equation}
\prod_{k=1}^{(n-1)/2} (-q^{n+2k};q^{2n})_\infty\,(-q^{n-2k};q^{2n})_\infty
= \prod_{\substack{r\ \mathrm{odd}\\ r \not\equiv n \pmod{2n}}} (1+q^r)
= \frac{\chi(q)}{\chi(q^n)},
\end{equation}
since $\chi(q) = \prod_{m\geq 0}(1+q^{2m+1})$ is the product over all odd exponents
and $\chi(q^n) = \prod_{m\geq 0}(1+q^{n(2m+1)})$ is the product over the odd multiples
of $n$, which are exactly the exponents congruent to $n$ modulo $2n$ (as $n$ is odd).

Therefore the numerator of $\prod_k u_k$ is
\begin{equation}
2^{(n-1)/2} q^{\sum k^2/n} \cdot \frac{\chi(q)}{\chi(q^n)}
\cdot (q^{2n};q^{2n})_\infty^{(n-1)/2},
\end{equation}
with $\sum_{k=1}^{(n-1)/2} k^2 = (n-1)n(n+1)/24$, so $\sum k^2/n = (n^2-1)/24$.
The denominator is $\varphi(q^n)^{(n-1)/2}$. By the triple product with $a=b=q^n$,
\begin{equation}
\varphi(q^n) = f(q^n,q^n) = (-q^n;q^{2n})_\infty^2\,(q^{2n};q^{2n})_\infty
= \chi(q^n)^2\,(q^{2n};q^{2n})_\infty,
\end{equation}
so the denominator is $\chi(q^n)^{n-1}(q^{2n};q^{2n})_\infty^{(n-1)/2}$. The
$(q^{2n};q^{2n})_\infty$ factors cancel, giving \eqref{eq:pprod-chi}.

For \eqref{eq:pprod-G}, substitute $\chi(q) = 2^{1/4}q^{1/24}G_n$ and
$\chi(q^n) = 2^{1/4}q^{n/24}G_{n^2}$:
\begin{equation}
\frac{\chi(q)}{\chi(q^n)^n}
= 2^{(1-n)/4} q^{(1-n^2)/24} \frac{G_n}{G_{n^2}^{\,n}},
\end{equation}
and multiplying by $2^{(n-1)/2}q^{(n^2-1)/24}$ gives
$2^{(n-1)/4}G_n/G_{n^2}^{\,n}$.
\end{proof}

\begin{corollary}
For $p=11$,
\begin{equation}
\pprod = 32\,q^5\,\frac{\chi(q)}{\chi^{11}(q^{11})}
= 4\sqrt{2}\,\frac{G_n}{G_{121n}^{11}}.
\end{equation}
At the flagship point $q = e^{-\pi/\sqrt{11}}$ this gives
\begin{equation}
\pprod = 4\sqrt{2}/G_{11}^{10} = 0.416439434275864095678569298496\dots.
\end{equation}
\end{corollary}

\begin{remark}
The $p=9$ (composite) case has a different shape: $u_1u_2u_4 =
8q^{7/3}\chi(q)/[\chi^6(q^9)\chi(q^3)]$ (Kim--Reb\'ak \cite{KR25}, Lemma 2.5),
because the middle term $u_3 = \varphi(q)/\varphi(q^9)-1$ is obtained separately
and the remaining terms are cube roots, not 9th roots. Theorem~\ref{thm:pprod}
applies to the product over all $(n-1)/2$ terms for odd $n$; for composite $n$ the
individual $u_k$ are not all $n$th roots of a single auxiliary polynomial, so the
product formula is a necessary but not sufficient ingredient for a composite-degree
theory.
\end{remark}

\section{Computational verification}
\label{sec:comp}

All computations in this section were performed in Python 3 with \texttt{sympy} 1.12
and \texttt{mpmath} 1.2.1, at 50 digits of precision unless otherwise stated. The
scripts are available from the author. Exact $q$-series were computed with integer
coefficients in the variable $t = q^{1/n}$, so all exponents are integers.

\subsection{The dissection}

We implemented the dissection \eqref{eq:dissection} in the variable $t = q^{1/n}$:
\begin{equation}
\frac{\varphi(t)}{\varphi(t^{n^2})}
= 1 + \sum_{k=1}^{(n-1)/2} \frac{2t^{k^2}\, f(t^{n(n+2k)},\, t^{n(n-2k)})}{\varphi(t^{n^2})},
\end{equation}
with exact integer-coefficient series for $\varphi$, $f$, and $\chi$ (series division
by inversion), plus high-precision numerics.

\begin{proposition}
For $n = 3, 5, 7, 9, 11, 13, 17$, the dissection holds exactly as a formal power
series to order $t^{120}$ (exact integer arithmetic), and numerically to 50 digits
with $|{\rm LHS} - {\rm RHS}| \leq 2\times 10^{-49}$ at $q = 0.3, 0.5, 0.7$.
The product formula \eqref{eq:pprod-chi} holds exactly to order $t^{120}$ and
numerically with relative error $\leq 6\times 10^{-50}$ for the same $n$.
\end{proposition}

\subsection{The three flagship evaluations}

The flagship evaluations are the payoff of the theory: exact values at CM points
where $\pprod$ and $Y$ become algebraic of low degree.
\begin{align}
\varphi(e^{-3\pi\sqrt{3}}) &= 3^{-3/4}\varphi(e^{-\pi\sqrt{3}})\{1+2^{1/3}\},
\label{eq:flag3} \\
\varphi(e^{-5\pi\sqrt{5}}) &= 5^{-3/4}\varphi(e^{-\pi\sqrt{5}})
\left\{1+\left(\frac{2(1+\tan\frac{\pi}{5})}{1-\sin\frac{\pi}{5}}\right)^{1/5}
+\left(\frac{2(1-\tan\frac{\pi}{5})}{1+\sin\frac{\pi}{5}}\right)^{1/5}\right\},
\label{eq:flag5} \\
\varphi(e^{-7\pi\sqrt{7}}) &= 7^{-3/4}\varphi(e^{-\pi\sqrt{7}})
\left\{1+\left(\frac{\cos\frac{\pi}{7}}{2\cos^2\frac{2\pi}{7}}\right)^{2/7}
+\left(\frac{\cos\frac{2\pi}{7}}{2\cos^2\frac{3\pi}{7}}\right)^{2/7}
+\left(\frac{\cos\frac{3\pi}{7}}{2\cos^2\frac{\pi}{7}}\right)^{2/7}\right\}.
\label{eq:flag7}
\end{align}

\begin{proposition}
The identities \eqref{eq:flag3}--\eqref{eq:flag7} hold numerically to 50 digits
(relative errors $0$ to $3\times 10^{-51}$). Moreover, the machine reproduction at
$q_1 = e^{-\pi/\sqrt{p}}$ confirms the internal structure of each theory:
$\pprod$ computed from the $\chi$-quotient equals $\pprod$ computed from the
$G$-formula; the multiplier relations hold; the auxiliary polynomials have exactly
the published coefficients; the roots match the trigonometric forms; and the cyclic
root ordering is as published.
\end{proposition}

In particular, for the septic case: $\pprod = 1$ (both routes, exact);
$X = \varphi^4(q)/\varphi^4(q^7) = 7$; the modular-equation residual is
$-2\times 10^{-49}$; $r(\xi) = \xi^3-6\xi^2+5\xi-1$ exactly; its roots match
$1/(2\cos k\pi/7)^2$; the cyclic assignment $(\alpha,\beta,\gamma)$ = roots for
$k=3,2,1$ gives $u = (\alpha^2\pprod/\beta)^{1/7}$ etc.\ matching the dissection's
$u,v,w$; and Lemma 3.8 of \cite{Reb23} is confirmed:
$\sum_i s_i^7 = 41.0$ (relative error $10^{-50}$), where $s_i$ are the published
trigonometric terms. We note for the record that the published septic terms appear
in the order $(v,u,w)$ relative to the dissection's $u,v,w$; the sum is unaffected.

For the quintic case: $\pprod = \sqrt{5}-1$ (both routes); $\varphi^2(q)/\varphi^2(q^5)
= \sqrt{5} = 1+\pprod$; the auxiliary quadratic $\xi^2-((\pprod-1)^2+7)\xi+\pprod^3=0$
has roots $\alpha \approx 6.7770610073$, $\beta \approx 0.2786670827$; and
$(\alpha\pprod)^{1/5}$, $(\beta\pprod)^{1/5}$ reproduce the two published
$\tan/\sin$ terms exactly.

\subsection{Class invariants}

We verified eleven class invariants against the definition
\begin{equation}
G_n = 2^{-1/4}q^{-1/24}\chi(q), \qquad q = e^{-\pi\sqrt{n}},
\end{equation}
to 50 digits:

\begin{center}
\small
\begin{tabular}{c|l|c}
$n$ & closed form & relative error \\ \hline
3 & $2^{1/12}$ & $3\times 10^{-51}$ \\
5 & $((1+\sqrt5)/2)^{1/4}$ & $2\times 10^{-51}$ \\
7 & $2^{1/4}$ & $0$ \\
11 & $[(3\sqrt{33}+17)^{1/3}-(3\sqrt{33}-17)^{1/3}+2]/(3\cdot 2^{1/4})$ & $0$ \\
25 & $(\sqrt5+1)/2$ & $7\times 10^{-51}$ \\
27 & $2^{1/12}(2^{1/3}-1)^{-1/3}$ & $3\times 10^{-51}$ \\
99 & $G_{99}^3 = \frac{2^{1/4}}{3}[(7822+1362\sqrt{33})^{1/3}$ & $2\times 10^{-51}$ \\
   & $\qquad +(7957+1383\sqrt{33})^{1/3}+2\sqrt{33}+13]$ & \\
729 & $G_{729}^3 = \frac{\sqrt2}{2}R^{1/3}
\left(\frac{3(\sqrt3+1)}{\sqrt3+1-R^{1/3}}-2\right)$ & $2\times 10^{-49}$ \\
   & $R = \frac{(2(\sqrt3+1))^{1/3}+1}{(2(\sqrt3-1))^{1/3}-1}$ & \\
125 & Berndt--Reb\'ak \cite{BR25}, Cor.~5.18 & $6\times 10^{-51}$ \\
625 & Berndt--Reb\'ak \cite{BR25}, Cor.~5.19 & $6\times 10^{-51}$ \\
243 & Watson: $2^{1/12}(2^{1/3}+2^{2/3}+3^{1/3})/(9-2\cdot 3^{4/3})^{1/3}$
& $2\times 10^{-51}$
\end{tabular}
\end{center}

Reciprocity $G_{1/n} = G_n$ was confirmed for $n=3,5,7,11$ (relative error
$\leq 7\times 10^{-51}$).

\begin{proposition}
The class invariant $G_{121}$ has the numerical value
\begin{equation}
G_{121} = 3.54880113413074478010727440546246876515786749179993826540053\dots
\end{equation}
to 60 digits. No closed form for $G_{121}$ is known in the literature. An integer-relation
search (PSLQ) found no minimal polynomial for $G_{121}$ (or $G_{121}^2$, $G_{121}^4$)
of degree $\leq 8$ at 80 digits.
\end{proposition}

\section{Sturm-bound proofs of the multiplier relations}
\label{sec:sturm}

The multiplier relations are the heart of each theory: they relate
$\varphi(q)/\varphi(q^p)$ to $\pprod$ and are what make the auxiliary polynomial
solvable. We prove them by the modern route: write every term as an eta-quotient,
verify modularity (Newman congruences, character, holomorphicity at all cusps via
Ligozat's formula), and check the $q$-expansion to the Sturm bound.

We use $\varphi(q) = \eta(2z)^5/(\eta(z)^2\eta(4z)^2)$ and
$\chi(q) = q^{1/24}\eta(2z)^2/(\eta(z)\eta(4z))$, where $\eta$ is the Dedekind eta
function and $q = e^{2\pi i z}$. For an eta-quotient $\prod_d \eta(dz)^{e_d}$ of
level $N$, the Newman congruences are $\sum_d d\,e_d \equiv 0$ and
$\sum_d (N/d)e_d \equiv 0 \pmod{24}$; the character is determined by
$D_{\rm eff} = (-1)^k\prod_d d^{e_d}$ (two terms share a character iff the ratio of
their $D_{\rm eff}$'s is a square in $\Q$); and the order at the cusp $1/c$ is given
by Ligozat's formula
$\operatorname{ord}_{1/c} = \frac{N}{24\gcd(c^2,N)}\sum_d e_d\gcd(c,d)^2/d$.
The Sturm bound is $B = \lceil k[\mathrm{SL}_2(\Z):\Gamma_0(N)]/12\rceil$ for weight $k$.

\begin{theorem}
The multiplier relations of degrees $3$, $5$, and $7$ hold as identities of modular
forms, proved by the Sturm-bound method:
\begin{center}
\begin{tabular}{c|c|c|c|c}
$p$ & identity & level & weight & Sturm bound \\ \hline
3 & $\varphi^4(q) - \varphi^4(q^3) - \pprod^3\varphi^4(q^3) = 0$ & $\Gamma_0(12)$ & 2 & 4 \\
5 & $\varphi^2(q) - \varphi^2(q^5) - \pprod\,\varphi^2(q^5) = 0$ & $\Gamma_0(20)$ & 1 & 3 \\
7 & $\varphi^8(q) - (2+5\pprod)\varphi^4(q)\varphi^4(q^7)
+ (1-\pprod)^3\varphi^8(q^7) = 0$ & $\Gamma_0(28)$ & 4 & 16
\end{tabular}
\end{center}
In each case all terms have the same weight and character (trivial for $p=3,7$;
$\chi_{-4}$ for all three $p=5$ terms), are holomorphic at every cusp, and the
$q$-expansion of the sum vanishes to order $q^{60}$, well beyond the Sturm bound.
\end{theorem}

\begin{proof}[Sketch of the verification]
For each identity we expand every term as an eta-quotient and check the Newman
congruences, the character (via $D_{\rm eff}$), and holomorphicity at every cusp
$1/c$, $c \mid N$, by Ligozat's formula. The $p=7$ identity expands to seven
eta-quotient terms after clearing $(1-\pprod)^3$ and $(2+5\pprod)$; all seven have
weight $4$, trivial character, and non-negative cusp orders. The $q$-expansion of
each sum is then verified to vanish through the Sturm bound (and in fact to $q^{60}$)
by exact integer arithmetic. The full eta-quotient data and the verification script
are available from the author.
\end{proof}

\begin{remark}
The degree-$11$ Schr\"oter equation (Berndt, Part III, Entry 7(i)),
$(\alpha\beta)^{1/4} + ((1-\alpha)(1-\beta))^{1/4}
+ 2(16\alpha\beta(1-\alpha)(1-\beta))^{1/12} = 1$ with
$\alpha = \lambda(\tau)$, $\beta = \lambda(11\tau)$, was verified numerically to
50 digits at $\tau = i, 0.5i, 2i, 0.3i$ ($|\mathrm{LHS}-1| \leq 2\times 10^{-47}$).
A full Sturm-bound proof was not carried out: the equation lives on
$\Gamma(2)\cap\Gamma_0(11)$, involves 4th and 12th roots, and clearing them to a
polynomial in $(\alpha,\beta)$ needs a resultant elimination that the present
eta-quotient machinery (integer-$d$ eta-quotients on $\Gamma_0(N)$ only) does not
support. Extending the prover to $\Gamma(2)$ is left for future work.
\end{remark}

\section{The degree-eleven modular equation}
\label{sec:p11eq}

We now prove the degree-$11$ modular equation: the explicit algebraic relation
between the multiplier $m = \varphi^2(q)/\varphi^2(q^{11})$ and the product
$\pprod = 32q^5\chi(q)/\chi^{11}(q^{11})$ (Theorem~\ref{thm:pprod} with $n=11$).
Only even powers of $m$ occur, so the relation is really in
$Y = m^2 = \varphi^4(q)/\varphi^4(q^{11})$, the exact analogue of the $p=7$
variable $X = \varphi^4(q)/\varphi^4(q^7)$. We write $P = \pprod/32$.

\begin{theorem}[The degree-$11$ modular equation]
\label{thm:p11}
Let $Y = \varphi^4(q)/\varphi^4(q^{11})$ and $P = \pprod/32$, where
$\pprod = 32q^5\chi(q)/\chi^{11}(q^{11})$. Then $G(Y,P) = 0$, where
\begin{align}
G(Y,P) ={}& -Y^5 + 5Y^4 - 10Y^3 + 10Y^2 - 5Y + 1 \nonumber \\
&+ P\,(160Y^4 + 2432Y^3 + 19392Y^2 + 10624Y + 160) \nonumber \\
&+ P^2\,(-10240Y^3 - 55296Y^2 - 796672Y + 10240) \nonumber \\
&+ P^3\,(327680Y^2 + 6914048Y + 327680) \nonumber \\
&+ P^4\,(-5242880Y + 5242880) + 33554432\,P^5.
\label{eq:G}
\end{align}
Equivalently, in the variables $(m,\pprod)$,
\begin{align}
F(m,\pprod) ={}& -m^{10} + 5m^8\pprod + 5m^8 - 10m^6\pprod^2 + 76m^6\pprod - 10m^6 \nonumber \\
&+ 10m^4\pprod^3 - 54m^4\pprod^2 + 606m^4\pprod + 10m^4 \nonumber \\
&- 5m^2\pprod^4 + 211m^2\pprod^3 - 778m^2\pprod^2 + 332m^2\pprod - 5m^2 \nonumber \\
&+ \pprod^5 + 5\pprod^4 + 10\pprod^3 + 10\pprod^2 + 5\pprod + 1 = 0.
\label{eq:F}
\end{align}
\end{theorem}

The relation has a striking structure. Writing out the pure parts,
\begin{equation}
G(Y,P) = (1-Y)^5 + (1+32P)^5 - 1 + P\cdot Y\cdot S(Y,P),
\label{eq:Gstruct}
\end{equation}
where $S$ is an explicit cubic, namely
\begin{align}
S(Y,P) = -32\,\big(&163840P^3 - 10240P^2Y - 216064P^2 + 320PY^2 + 1728PY \nonumber \\
&+ 24896P - 5Y^3 - 76Y^2 - 606Y - 332\big).
\end{align}
The $P$-only part is $(1+32P)^5 = (1+\pprod)^5$ (coefficients $1, 160, 10240,
327680, 5242880, 33554432 = \binom{5}{k}2^{5k}$) and the $Y$-only part is
$(1-Y)^5$. Only monomials $Y^iP^j$ with $i+j \leq 5$ occur, and $G$ is irreducible
over $\Q$.

\subsection{Where $m$ and $\pprod$ live}

Both $m$ and $\pprod$ are eta-quotients on $\Gamma_0(44)$ --- level $4\cdot 11$,
exactly as in the $p=7$ case (level $28$), \emph{not} level $11$:
\begin{align}
m &= \eta(z)^{-4}\eta(2z)^{10}\eta(4z)^{-4}\eta(11z)^4\eta(22z)^{-10}\eta(44z)^4, \\
\pprod &= 32\cdot\eta(z)^{-1}\eta(2z)^2\eta(4z)^{-1}\eta(11z)^{11}\eta(22z)^{-22}\eta(44z)^{11}.
\end{align}
The eta-quotient for $\pprod$ itself has $q$-order $5$ at $\infty$, so
$\pprod = 32q^5\cdot(\text{unit})$, as in Theorem~\ref{thm:pprod}.
Both satisfy the Newman congruences at level $44$ with trivial character
($D_{\rm eff}(m) = 1/121$ and $D_{\rm eff}(\pprod) = 1$, both squares), so
$m, \pprod \in \C(X_0(44))$. Their Ligozat cusp orders on $\Gamma_0(44)$ are:
\begin{center}
\begin{tabular}{c|c|c}
cusp & $\operatorname{ord}(m)$ & $\operatorname{ord}(\pprod)$ \\ \hline
$\infty$ & $0$ (value $1$) & $5$ (zero) \\
$0$ & $0$ (value $11$) & $0$ (value $32$) \\
$1/2$ & $5$ (zero) & $0$ \\
$1/4$ & $0$ & $0$ \\
$1/11$ & $0$ & $5$ (zero) \\
$1/22$ & $-5$ (pole) & $-10$ (pole)
\end{tabular}
\end{center}
Hence $\deg(m) = 5$, $\deg(\pprod) = 10$, and both functions have their only pole
at the cusp $1/22$. The degree bounds for the minimal relation are therefore
$\deg_m F \leq 10$ and $\deg_{\pprod} F \leq 5$. The cusp values check out:
$F(1,0) = 0$ and $F(11,32) = 0$, equivalently $G(1,0) = G(121,1) = 0$.

\subsection{How the relation was found}

Exact integer $q$-series for $m$ and $P = \pprod/32$ were computed to $q^{400}$.
Monomials $m^iP^j$ were built by convolution modulo the prime $999983$ (float64 is
exact here: partial sums $< 6\cdot 10^{14} < 2^{53}$), and the nullspace was computed
over increasing monomial boxes $(I,J)$. No relation exists for any box below
$(10,5)$; at $(10,5)$ (66 monomials, 106 rows) the nullspace is one-dimensional,
supported only on even powers of $m$ with the triangular pattern $Y^iP^j$,
$i+j \leq 5$, and $P^5$-coefficient $1$. Exact rational nullspace computation on the
21 triangle monomials then gives the unique primitive integer vector --- the
coefficients in \eqref{eq:G} (largest $33554432 = 2^{25}$; all coefficients are
small). No PSLQ or LLL was needed: exact linear algebra sufficed.

\subsection{Proof by the Sturm bound}

\begin{proof}[Proof of Theorem~\ref{thm:p11}]
On $\Gamma_0(44)$, the index is $\mu = [\mathrm{SL}_2(\Z):\Gamma_0(44)] = 72$.
Let $K = \eta(2z)^2\eta(22z)^2$ (weight $2$, trivial character). Then $K^{25}$ has
weight $50$ and $\operatorname{ord}_{1/22}(K^{25}) = 50$, while $Y^iP^j$ has its
only pole at $1/22$, of order $10(i+j) \leq 50$. Hence every term
$K^{25}Y^iP^j$ (21 terms) is a holomorphic modular form of weight $50$, trivial
character, on $\Gamma_0(44)$ --- each term's Newman congruences, $D_{\rm eff}$
square, and Ligozat orders $\geq 0$ at every cusp were checked symbolically.
The Sturm bound is $\lceil 50\cdot 72/12\rceil = 300$. The $q$-expansion of
$K^{25}G(Y,P)$, computed in exact integer arithmetic, has all coefficients
$q^0,\dots,q^{300}$ equal to zero (in fact no nonzero coefficient to $q^{400}$).
Therefore $K^{25}G(Y,P) \equiv 0$ as a modular form, and since $K \not\equiv 0$,
$G(Y,P) = 0$ identically. This is a complete proof in the same style as the
proofs in \S\ref{sec:sturm}.
\end{proof}

The relation was also verified independently: as a formal power series
$G(Y,P) = 0$ to $q^{600}$ in exact integer arithmetic, and numerically to 50
digits with relative error $\leq 7\times 10^{-50}$ at eight real $q \in
[0.2, 0.85]$, at the flagship $q_1 = e^{-\pi/\sqrt{11}}$ (where $m = \sqrt{11}$
and $\pprod = 4\sqrt{2}/G_{11}^{10} = 0.41643943427586409568\dots$), and at three
complex $q$.

\subsection{The curve $G(Y,P) = 0$ has genus $2$}

The curve is irreducible over $\Q$ and is a plane quintic. Its degree-$5$
homogeneous part is $-(Y-32P)^5$, so the projective curve meets the line at
infinity only at $[32:1:0]$, which is smooth (the line at infinity is tangent
with contact order $5$). It has exactly four singular points in the affine plane,
all ordinary double points (nodes, $\delta = 1$ each), whose $Y$-coordinates are
the four roots of $59049Y^4 + 280422Y^3 + 293768Y^2 - 64614Y - 14641$
($59049 = 3^{10}$, $14641 = 11^4$). Hence the genus is
$(5-1)(5-2)/2 - \sum\delta = 6 - 4 = 2$.

\begin{remark}
This corrects an expectation in earlier work on this problem, which predicted a
genus-$1$ (elliptic) relation on the grounds that $X_0(11)$ has genus $1$. The
functions $m$ and $\pprod$ actually live on $X_0(44)$, not $X_0(11)$, and their
image curve has genus $2$. The map $X_0(44) \to C$ has degree $2$ (the tower
$[\C(X_0(44)):\Q(P)] = 10$, $[\Q(Y,P):\Q(P)] = 5$), and Riemann--Hurwitz is
consistent: $g(X_0(44)) = 4$, $g(C) = 2$, so $2\cdot 4 - 2 = 2(2\cdot 2 - 2) + R$
gives $R = 2$ ramification points. The degree-$2$ map is the Atkin--Lehner
quotient by $w_4$: the genuine involution $w_4$ of $X_0(44)$ fixes both $Y$ and
$P$ (verified numerically to 40 digits), so $C \cong X_0(44)/w_4$ and the field
extension is $\Q(m,P)/\Q(Y,P)$ with Galois involution $m \mapsto -m$. The
involution $w_{11}$ (equivalently $w_{44}$) sends $Y \mapsto 121/Y$ and preserves
the curve, but no simple closed form for its action on $P$ was found.
\end{remark}

\section{The degree-thirteen and degree-seventeen modular equations}
\label{sec:p1317}

The method of \S\ref{sec:p11eq} works prime by prime: for $p = 13$ and $p = 17$
the functions $m$ and $\pprod$ are again eta-quotients on $\Gamma_0(4p)$ with
trivial character and a single pole at the cusp $1/(2p)$, so the same search
(mod-prime nullspace on $q$-expansions) and the same proof (Sturm bound after
clearing the pole with a power of $K = \eta(2z)^2\eta(2pz)^2$) apply. The two
new cases differ from $p = 11$ in one respect: odd powers of $m$ occur, so the
relations are in $(m,P)$, not in $(Y,P)$.

\subsection{Degree $13$}

Here $\pprod = 64q^7\chi(q)/\chi^{13}(q^{13})$ (Theorem~\ref{thm:pprod} with
$n = 13$) and $P = \pprod/64$. The eta-quotients on $\Gamma_0(52)$ are
\begin{align}
m &= \eta(z)^{-4}\eta(2z)^{10}\eta(4z)^{-4}\eta(13z)^4\eta(26z)^{-10}\eta(52z)^4, \\
\pprod &= 64\cdot\eta(z)^{-1}\eta(2z)^2\eta(4z)^{-1}\eta(13z)^{13}
\eta(26z)^{-26}\eta(52z)^{13}.
\end{align}
The eta-quotient for $\pprod$ itself has $q$-order $7$ at $\infty$, so
$\pprod = 64q^7\cdot(\text{unit})$, as in Theorem~\ref{thm:pprod}.
Both satisfy the Newman congruences at level $52$ with trivial character
($D_{\rm eff}(m) = 1/169$ and $D_{\rm eff}(P) = 1$, both squares). Their
Ligozat cusp orders are: $m$ has order $6$ at $1/2$, $-6$ at $1/26$, and $0$
elsewhere; $P$ has order $7$ at $\infty$ and $1/13$, $-14$ at $1/26$,
and $0$ elsewhere. (The cusp $1/52$ is $\Gamma_0(52)$-equivalent to $\infty$.)
Hence $\deg(m) = 6$, $\deg(P) = 14$, and both have their
only pole at $1/26$.

\begin{theorem}[The degree-$13$ modular equation]
\label{thm:p13}
Let $m = \varphi^2(q)/\varphi^2(q^{13})$ and $P = \pprod/64$, where
$\pprod = 64q^7\chi(q)/\chi^{13}(q^{13})$. Then $F(m,P) = 0$, where
\begin{align}
F(m,P) ={}& 262144P^3 + 45056P^2m^2 - 188416P^2m + 12288P^2 \nonumber \\
&+ 1216Pm^4 - 2816Pm^3 + 10368Pm^2 + 7424Pm + 192P \nonumber \\
&- m^7 + 7m^6 - 21m^5 + 35m^4 - 35m^3 + 21m^2 - 7m + 1.
\label{eq:F13}
\end{align}
The pure parts are $F(m,0) = -(m-1)^7$ and $F(0,P) = (1+64P)^3$, and
\begin{equation}
F = (1+64P)^3 - (m-1)^7 - 1
+ 64\cdot P\cdot m\cdot(704Pm - 2944P + 19m^3 - 44m^2 + 162m + 116).
\label{eq:F13struct}
\end{equation}
$F$ is irreducible over $\Q$.
\end{theorem}

\begin{proof}[Proof of Theorem~\ref{thm:p13}]
The relation first appears at the monomial box $(7,4)$: no relation exists
for any smaller box, and at $(7,4)$ the nullspace is two-dimensional (the
relation and its multiple by $P$), supported on 17 monomials
($m^0,\dots,m^7$; $m^0,\dots,m^4$ times $P$; $m^0,m^1,m^2$ times $P^2$;
$P^3$). Exact rational nullspace computation gives the unique primitive
integer vector --- the coefficients in \eqref{eq:F13}.

For the proof, $[\mathrm{SL}_2(\Z):\Gamma_0(52)] = 84$. Let
$K = \eta(2z)^2\eta(26z)^2$ (weight $2$, trivial character); $K$ has order
$7/3$ at every cusp. Then $K^{18}$ has weight $36$ and
$\operatorname{ord}_{1/26}(K^{18}) = 42$, equal to the maximum pole order of
$m^iP^j$ over the support. Hence every term $K^{18}m^iP^j$ (17 terms) is a
holomorphic modular form of weight $36$, trivial character, on $\Gamma_0(52)$
--- each term's Newman congruences, $D_{\rm eff}$ square, and Ligozat orders
$\geq 0$ at every cusp were checked symbolically. The Sturm bound is
$\lceil 36\cdot 84/12\rceil = 252$. The $q$-expansion of $K^{18}F(m,P)$,
computed in exact integer arithmetic, has all coefficients $q^0,\dots,q^{252}$
equal to zero (in fact no nonzero coefficient to $q^{400}$). Therefore
$K^{18}F(m,P) \equiv 0$ as a modular form, and since $K \not\equiv 0$,
$F(m,P) = 0$ identically.
\end{proof}

The relation was also verified independently: as a formal power series
$F(m,P) = 0$ to $q^{600}$ in exact integer arithmetic, and numerically to 50
digits with relative error $\sim 10^{-51}$ at seven points, including the
flagship $q_1 = e^{-\pi/\sqrt{13}}$ (where $m = \sqrt{13}$ and
$P = 0.0034695471649933081995\dots$). The cusp values check out:
$F(1,0) = 0$ and $F(13,1) = 0$.

The curve $F(m,P) = 0$ is a plane septic of arithmetic genus $15$. It has one
point at infinity, unibranch with Puiseux pair $(4,7)$ and $\delta = 9$, and
exactly four affine nodes ($\delta = 1$ each) at $(1,1/4)$,
$(33.54542973059973, -93.17083713060966)$,
$(-2.382790973655804, -0.086883662251872)$, and
$(-0.162638756943925, 0.007720792861530)$. Hence its genus is $15 - 9 - 4 = 2$,
confirmed by Riemann--Hurwitz on $m$ ($R = 8$: $2$ at $m = 0$, $4$ at the four
simple ramification points, $2$ at $m = \infty$; the nodes contribute $0$) and
by the degree-$2$ map $X_0(52) \to C$ ($g(X_0(52)) = 5$, $R' = 4$). The fibre
over $m = 0$ is the single point $(0,-1/64)$, smooth, with $e = 3$
($F_m(0,-1/64) = -169 = -13^2 \ne 0$).

\subsection{Degree $17$}

Here $\pprod = 256q^{12}\chi(q)/\chi^{17}(q^{17})$ (Theorem~\ref{thm:pprod}
with $n = 17$) and $P = \pprod/256$. The eta-quotients on $\Gamma_0(68)$ are
\begin{align}
m &= \eta(z)^{-4}\eta(2z)^{10}\eta(4z)^{-4}\eta(17z)^4\eta(34z)^{-10}\eta(68z)^4, \\
\pprod &= 256\cdot\eta(z)^{-1}\eta(2z)^2\eta(4z)^{-1}\eta(17z)^{17}
\eta(34z)^{-34}\eta(68z)^{17}.
\end{align}
The eta-quotient for $\pprod$ itself has $q$-order $12$ at $\infty$, so
$\pprod = 256q^{12}\cdot(\text{unit})$, as in Theorem~\ref{thm:pprod}.
Both satisfy the Newman congruences at level $68$ with trivial character
($D_{\rm eff}(m) = 1/289$ and $D_{\rm eff}(P) = 1$, both squares). Their
Ligozat cusp orders are: $m$ has order $8$ at $1/2$, $-8$ at $1/34$, and $0$
elsewhere; $P$ has order $12$ at $\infty$ and $1/17$, $-24$ at $1/34$,
and $0$ elsewhere. (The cusp $1/68$ is $\Gamma_0(68)$-equivalent to $\infty$.)
Hence $\deg(m) = 8$, $\deg(P) = 24$, and both have their
only pole at $1/34$.

\begin{theorem}[The degree-$17$ modular equation]
\label{thm:p17}
Let $m = \varphi^2(q)/\varphi^2(q^{17})$ and $P = \pprod/256$, where
$\pprod = 256q^{12}\chi(q)/\chi^{17}(q^{17})$. Then $F(m,P) = 0$, where
$F$ has the $35$ monomials with $i + 3j \leq 12$:
\begin{align}
F(m,P) ={}& \phantom{+} m^{12} - 12m^{11} + 66m^{10} - 220m^9 + 495m^8
- 792m^7 + 924m^6 - 792m^5 + 495m^4 \nonumber \\
&- 220m^3 + 66m^2 - 12m + 1 \nonumber \\
&+ P\,(-1024 - 391168m - 5329920m^2 - 10408960m^3 + 5120m^4 + 19456m^5
\nonumber \\
&\qquad - 744448m^6 + 99328m^7 - 24576m^8 - 1024m^9) \nonumber \\
&+ P^2\,(393216 - 195100672m + 919994368m^2 - 192806912m^3
+ 107347968m^4 \nonumber \\
&\qquad - 19464192m^5 + 393216m^6) \nonumber \\
&+ P^3\,(-67108864 - 6643777536m - 3690987520m^2 - 67108864m^3)
+ 4294967296P^4.
\label{eq:F17}
\end{align}
The pure parts are $F(m,0) = (m-1)^{12}$ and $F(0,P) = (256P-1)^4$, and
\begin{equation}
F = (m-1)^{12} + (256P-1)^4 - 1 - 1024\cdot P\cdot m\cdot S(m,P),
\label{eq:F17struct}
\end{equation}
where $S$ is explicit:
\begin{align}
S(m,P) ={}& 65536P^2m^2 + 3604480P^2m + 6488064P^2
- 384Pm^5 + 19008Pm^4 - 104832Pm^3 \nonumber \\
&+ 188288Pm^2 - 898432Pm + 190528P \nonumber \\
&+ m^8 + 24m^7 - 97m^6 + 727m^5 - 19m^4 - 5m^3 + 10165m^2 + 5205m + 382.
\end{align}
$F$ is irreducible over $\Q$.
\end{theorem}

\begin{proof}[Proof of Theorem~\ref{thm:p17}]
The relation first appears at the monomial box $(12,4)$: no relation exists
for any smaller box, and at $(12,4)$ the nullspace is one-dimensional,
supported on the $35$ monomials with $i + 3j \leq 12$. Exact rational
nullspace computation gives the unique primitive integer vector --- the
coefficients in \eqref{eq:F17}.

For the proof, $[\mathrm{SL}_2(\Z):\Gamma_0(68)] = 108$. Let
$K = \eta(2z)^2\eta(34z)^2$ (weight $2$, trivial character); $K$ has order $3$
at every cusp. Then $K^{32}$ has weight $64$ and
$\operatorname{ord}_{1/34}(K^{32}) = 96$, equal to the maximum pole order of
$m^iP^j$ over the support (attained uniformly on the triangle boundary
$i + 3j = 12$). Hence every term $K^{32}m^iP^j$ (35 terms) is a holomorphic
modular form of weight $64$, trivial character, on $\Gamma_0(68)$ --- each
term's Newman congruences, $D_{\rm eff}$ square, and Ligozat orders $\geq 0$
at every cusp were checked symbolically. The Sturm bound is
$\lceil 64\cdot 108/12\rceil = 576$. The $q$-expansion of $K^{32}F(m,P)$,
computed in exact integer arithmetic, has all coefficients $q^0,\dots,q^{576}$
equal to zero (in fact no nonzero coefficient to $q^{800}$). Therefore
$K^{32}F(m,P) \equiv 0$ as a modular form, and since $K \not\equiv 0$,
$F(m,P) = 0$ identically.
\end{proof}

The relation was also verified independently: as a formal power series
$F(m,P) = 0$ to $q^{600}$ in exact integer arithmetic, and numerically to 50
digits with relative error $\sim 10^{-51}$ at seven points, including the
flagship $q_1 = e^{-\pi/\sqrt{17}}$ (where $m = \sqrt{17}$ and
$P = 0.00017769354832319060602\dots$). The cusp values check out:
$F(1,0) = 0$ and $F(17,1) = 0$.

The curve $F(m,P) = 0$ is a plane curve of degree $12$ and arithmetic genus
$55$. It has one point at infinity, unibranch, with multiplicity sequence
$8, 4, 4$ under iterated blow-up, so $\delta = 28 + 6 + 6 = 40$ there. It has
exactly twelve affine nodes ($\delta = 1$ each); their $m$-coordinates are the
twelve roots of
\begin{align}
&15625m^{12} + 57350m^{11} + 435636m^{10} + 4017474m^9 - 7136763m^8
+ 58172316m^7 \nonumber \\
&\qquad - 2075760m^6 + 39725924m^5 + 142214563m^4
+ 63554366m^3 + 7265596m^2 - 57222m - 4913,
\end{align}
each verified a node by a nonzero tangent-cone discriminant
$F_{mm}F_{PP} - F_{mP}^2$ at 80 digits. Hence its genus is $55 - 40 - 12 = 3$,
confirmed by Riemann--Hurwitz on $m$ ($R = 12$: $3$ at $m = 0$, $6$ at the six
simple ramification points, $3$ at $m = \infty$; the nodes contribute $0$) and
by the degree-$2$ map $X_0(68) \to C$ ($g(X_0(68)) = 7$, $R' = 4$). The fibre
over $m = 0$ is the single point $(0,1/256)$, smooth, with $e = 4$.

\begin{remark}
Ramanujan already recorded degree-$17$ modular equations (Notebook 2,
Entry 12) --- but in the classical $(\alpha,\beta)$ variables, not in the
theta variables $(m,\pprod)$. Theorem~\ref{thm:p17} is the first degree-$17$
equation in Ramanujan's theta variables. Note also that $17$ is in Ramanujan's
``quadratics only'' list, since $(17-1)/2 = 8 = 2^3$, while $13$ is in his
``cubics'' list; the auxiliary polynomials $r_{13}$ and $r_{17}$ (degrees $6$
and $8$) remain open, exactly as $r_{11}$ is open.
\end{remark}

\subsection{The emerging pattern, and its limits}

\begin{center}
\begin{tabular}{c|c|c|c|c|c|c|c}
$p$ & level & index & relation & support & weight & Sturm bound & genus \\ \hline
11 & 44 & 72 & $G(Y,P)$, deg $5$ & $i+j \leq 5$ (21 mon.) & 50 & 300 & 2 \\
13 & 52 & 84 & $F(m,P)$, deg $7$ & 17 mon. & 36 & 252 & 2 \\
17 & 68 & 108 & $F(m,P)$, deg $12$ & $i+3j \leq 12$ (35 mon.) & 64 & 576 & 3
\end{tabular}
\end{center}

The common features are real: the pure parts are perfect powers with exponents
equal to the degrees ($\deg_m F$ and $\deg_P F$); the mixed remainder is
divisible by $P\cdot m$; and the map $X_0(4p) \to C$ has degree $2$ in all three
cases. But a uniform theorem for all primes $p$ is not established: the relation
is in $(m,P)$ for $p = 13, 17$ but in $(Y,P) = (m^2,P)$ for $p = 11$, with no
proved rule for when odd powers of $m$ appear; the support pattern differs in
each case; and the genus sequence $2, 2, 3$ has no proved formula. The method is
effective prime by prime --- the Sturm bound and the linear algebra grow with
$p$, but both scaled comfortably (seconds) at $p = 17$ --- and the next cases
are $p = 19$ (level $76$) and $p = 23$ (level $92$).

\section{A conjectural template for general prime $p$}
\label{sec:conjecture}

The six proved cases $p = 3, 5, 7, 11, 13, 17$ suggest a uniform template for
the degree-$p$ modular equation. We state it here precisely, and we label each
part honestly: \emph{[Theorem]} for what is proved in general, \emph{[Verified]}
for what has been checked numerically but not proved, \emph{[Pattern]} for what
holds in all cases where it can be checked but has no proof and no known
mechanism, and \emph{[Guess]} for extrapolation with thin evidence. Nothing
labelled [Pattern] or [Guess] is presented as a theorem. The prime $p = 3$ is
degenerate in every formula below (since $(p^2-1)/24$ is not an integer there)
and is excluded from the general statements.

Let $p \geq 5$ be prime, and write $m = \varphi^2(q)/\varphi^2(q^p)$ and
$P = \pprod/2^{(p-1)/2}$, where $\pprod$ is as in Theorem~\ref{thm:pprod}.
Both are eta-quotients on $\Gamma_0(4p)$:
\begin{align}
m &= \eta(z)^{-4}\eta(2z)^{10}\eta(4z)^{-4}\eta(pz)^4\eta(2pz)^{-10}\eta(4pz)^4, \\
P &= \eta(z)^{-1}\eta(2z)^2\eta(4z)^{-1}\eta(pz)^p\eta(2pz)^{-2p}\eta(4pz)^p.
\end{align}
A direct Ligozat computation at the six cusps $\infty$, $0$, $1/2$, $1/4$,
$1/p$, $1/(2p)$ of $\Gamma_0(4p)$ gives [Theorem]:
\begin{center}
\begin{tabular}{c|c|c}
cusp & $\operatorname{ord}(m)$ & $\operatorname{ord}(P)$ \\ \hline
$\infty$ & $0$ & $(p^2-1)/24$ \\
$0$ & $0$ & $0$ \\
$1/2$ & $(p-1)/2$ & $0$ \\
$1/4$ & $0$ & $0$ \\
$1/p$ & $0$ & $(p^2-1)/24$ \\
$1/(2p)$ & $-(p-1)/2$ & $-(p^2-1)/12$
\end{tabular}
\end{center}
Hence $\deg(m) = (p-1)/2$ and $\deg(P) = (p^2-1)/12$; each has its only pole
at $1/(2p)$, and $P$ has its only zeros at $\infty$ and $1/p$. The genus of
$X_0(4p)$ is $(p-3)/2$.

The key mechanism is the Atkin--Lehner involution $w_4$ of $X_0(4p)$, attached
to $Q = 4$. [Verified] Numerically --- eta via the pentagonal series, at 60
digits and three test points, for each $p$ in $\{3, 5, 7, 11, 13, 17\}$ ---
\begin{equation}
w_4(m) = \left(\frac{-1}{p}\right)\cdot m, \qquad w_4(P) = P
\qquad (p \geq 5),
\label{eq:w4}
\end{equation}
to within $10^{-45}$ (for $p = 3$, $w_4(m) = -m$ but $w_4(P) = \zeta_3 P$).
A proof through the eta-multiplier (Dedekind sums) for the representative of
$w_4$ has not been written; it is the single most important missing proof in
this template, because everything that follows depends on it. From
\eqref{eq:w4} we obtain the rule for which variables the relation uses
[Theorem, given \eqref{eq:w4}]: if $p \equiv 1 \pmod 4$ then $(-1/p) = +1$,
$w_4$ fixes $m$, and the relation is in $(m,P)$; if $p \equiv 3 \pmod 4$ then
$(-1/p) = -1$, $w_4$ sends $m \mapsto -m$, so only $Y = m^2$ is
$w_4$-invariant and the relation is in $(Y,P)$. This fits all six cases
($p = 3, 7, 11$ in $Y$; $p = 5, 13, 17$ in $m$).

\begin{conjecture}[The degree-$p$ modular equation]
\label{conj:template}
For every prime $p \geq 5$, with $X = m$ if $p \equiv 1 \pmod 4$ and
$X = Y = m^2$ if $p \equiv 3 \pmod 4$, there is an irreducible polynomial
$F_p(X,P) \in \Z[X,P]$ with $F_p(X,P) = 0$, of the form
\begin{equation}
F_p(X,P) = (1-X)^{e_X} + \left(1 - \left(\frac{2}{p}\right)2^{(p-1)/2}P\right)^{e_P}
- 1 + P\cdot X\cdot S_p(X,P),
\label{eq:template}
\end{equation}
where $e_X = (p^2-1)/24$ and $e_P = (p-1)/4$ if $p \equiv 1 \pmod 4$,
$e_P = (p-1)/2$ if $p \equiv 3 \pmod 4$, for some polynomial $S_p$; the
support of $F_p$ is the full triangle
$\{(i,j) : i\,e_P + j\,e_X \leq e_X e_P\}$; the curve $F_p = 0$ is the quotient
$X_0(4p)/w_4$; and its genus is $(p-1-R')/4$, where $R' = \#\operatorname{Fix}(w_4)$.
\end{conjecture}

The parts of Conjecture~\ref{conj:template} have very different status, and we
separate them.

\emph{The exponents [Theorem, given \eqref{eq:w4} and the degree-$2$ map].}
Since $\deg_{X_0(4p)}(P) = (p^2-1)/12$ and the map $X_0(4p) \to C$ has degree
$2$ (proved for $p \leq 17$; see below), $\deg_C(P) = (p^2-1)/24$, which equals
$e_X = [\,\Q(X,P):\Q(P)\,]$. Similarly $e_P = \deg_C(X)$. The $p = 13$
``anomaly'' ($e_X = 7 > e_P = 3$) is no anomaly: $e_X$ grows like $p^2/24$
while $e_P$ grows like $p/4$ or $p/2$, so $e_X > e_P$ for all $p \geq 11$,
with crossover exactly at $p = 11$.

\emph{The pure parts [Pattern, with a proved mechanism for the $X$-part].}
The template \eqref{eq:template} fits all six cases, including the genus-$0$
ones. The $X$-part is a perfect power by a proved mechanism, given
\eqref{eq:w4}: $P$'s zeros on $C$ are the single point $(X,P) = (1,0)$ --- the
image of the two cusps $\infty$ and $1/p$, which $w_4$ swaps, and at both of
which $X = 1$ (since $m(\infty) = 1$ and $m(1/p) = (-1/p)$, the latter
following from \eqref{eq:w4} and $w_4$ swapping $\infty \leftrightarrow 1/p$).
Since $\deg_C(P) = e_X$, the fibre $P = 0$ is the single point $X = 1$ with
multiplicity $e_X$. The $P$-part is a perfect power because $X$'s only zero on
$X_0(4p)$ is the cusp $1/2$, which is fixed by $w_4$ (hence ramified in the
degree-$2$ map), so $X$ has a single zero on $C$, of order $e_P$, at the
$P$-value $P(1/2)$. The sign is $-(2/p)$: the cusp value
$P(1/2) = (2/p)/2^{(p-1)/2}$ was verified numerically for $p = 5, 7, 11, 13, 17$
and for $p = 19, 23$ (relative error $10^{-11}$--$10^{-24}$ at
$\tau = 1/2 + 0.005i$), and the cusp values $m(1/p) = (-1/p)$ for
$p = 11, 13, 17, 19, 23$ follow from the deduction above. No proof of the
$(2/p)$ sign is known; the mechanism should be the eta-quotient cusp-value
formula at $1/2$.

\emph{The support [Pattern, 6/6].} The support of $F_p$ is the full triangle
with vertices $(0,0)$, $(e_X,0)$, $(0,e_P)$, every monomial occurring with
nonzero coefficient. This holds for $p = 11$ ($i+j \leq 5$, 21 monomials),
$p = 13$ ($3i+7j \leq 21$, 17), $p = 17$ ($i+3j \leq 12$, 35), and the three
genus-$0$ cases. No proof is known; it is a genericity statement, and the most
``accidental''-looking part of the template. The mixed remainder's divisibility
by $P\cdot X$ is then trivial algebra [Theorem, given the pure parts]: it
vanishes at $X = 0$ and at $P = 0$, hence is divisible by both.

\emph{The degree-$2$ map [Proved for $p \leq 17$; conjectured in general].}
$[\C(X_0(4p)):\Q(P)] = (p^2-1)/12$ and $[\Q(X,P):\Q(P)] = e_X = (p^2-1)/24$,
so the map $X_0(4p) \to C$ has degree $2$ in every proved case, and
$\Q(X,P) = \operatorname{Fix}(w_4)$ since both are index-$2$ subfields. The
general case needs only the existence of a relation $F_p$ with
$\deg_X F_p = (p^2-1)/24$, which is what the Sturm computation produces.
Irreducibility follows: $F_p$ is the minimal polynomial of $X$ over $\Q(P)$.

\emph{The genus [Guess, and the weakest part].} By Riemann--Hurwitz for the
degree-$2$ map, $g(C) = (p-1-R')/4$ with $R' = \#\operatorname{Fix}(w_4)$.
Two fixed points are cusps ($1/2$ and $1/(2p)$ are fixed; $w_4$ swaps
$\infty \leftrightarrow 1/p$ and $0 \leftrightarrow 1/4$), so
$R' = 2 + N_{\mathrm{CM}}(p)$ with $N_{\mathrm{CM}}(p)$ the number of fixed CM
points. The data are
\begin{center}
\begin{tabular}{c|c|c|c|c|c|c}
$p$ & 3 & 5 & 7 & 11 & 13 & 17 \\ \hline
$R'$ & 2 & 4 & 6 & 2 & 4 & 4 \\
$N_{\mathrm{CM}}(p)$ & 0 & 2 & 4 & 0 & 2 & 2 \\
$g(C)$ & 0 & 0 & 0 & 2 & 2 & 3
\end{tabular}
\end{center}
No formula for $N_{\mathrm{CM}}(p)$ was found: it is not a function of $p \bmod 4$
($p = 7$ and $p = 11$ are both $\equiv 3 \pmod 4$ but have $N_{\mathrm{CM}} = 4$
and $0$), nor of $p \bmod 8$. The count should come from the CM theory of $w_4$
(class numbers of orders in $\Q(i)$ and the number of $w_4$-admissible level
structures), but the derivation was not completed. The minimal-$R'$ prediction
($R' = 2$ if $p \equiv 3 \pmod 4$, $R' = 4$ if $p \equiv 1 \pmod 4$) fits
$p = 11, 13, 17$ but fails $p = 7$ (which has $R' = 6$, $g = 0$ instead of the
predicted $g = 1$); the genus predictions below are therefore the least
reliable part of this section.

\subsection{Predictions for $p = 19$ and $p = 23$}

\begin{center}
\begin{tabular}{c|c|c}
& $p = 19$ & $p = 23$ \\ \hline
variables $X$ & $Y = m^2$ & $Y = m^2$ \\
$e_X = (p^2-1)/24$ & $15$ & $22$ \\
$e_P = (p-1)/2$ & $9$ & $11$ \\
pure $X$-part & $(1-Y)^{15}$ & $(1-Y)^{22}$ \\
$(2/p)$ & $-1$ ($19 \equiv 3 \bmod 8$) & $+1$ ($23 \equiv 7 \bmod 8$) \\
pure $P$-part & $(1+512P)^9$ & $(1-2048P)^{11}$ \\
support & $3i+5j \leq 45$ & $i+2j \leq 22$ \\
monomials & $82$ & $144$ \\
level / index & $\Gamma_0(76)$ / $120$ & $\Gamma_0(92)$ / $144$ \\
genus $g(C)$, minimal-$R'$ guess & $4$ & $5$
\end{tabular}
\end{center}

The cusp values needed for the pure parts ($m(1/p) = -1$,
$P(1/2) = -1/512$ for $p = 19$; $P(1/2) = +1/2048$ for $p = 23$) are already
verified numerically, as above. The test procedure is the same Sturm pipeline
used for $p = 11, 13, 17$ (\S\ref{sec:p11eq}, \S\ref{sec:p1317}), with
$K = \eta(2z)^2\eta(2pz)^2$ raised to a power clearing the maximal pole order
over the predicted triangle. Each part of Conjecture~\ref{conj:template} has
its own falsifier: a $w_4$-ratio that is not the constant $(-1/p)$ (for $m$) or
$1$ (for $P$) kills the variable rule and the degree-$2$ identification; an
exponent or pure part different from the prediction kills \eqref{eq:template};
a missing interior monomial kills the support triangle (the most likely part
to fail); a genus larger than the minimal-$R'$ prediction would falsify the
framework itself, since $R' \geq 2$ always.

\section{The $p=11$ frontier}
\label{sec:p11}

\subsection{What the modular equation does and does not give}

Theorem~\ref{thm:p11} supplies Layer 3 of the uniform template for $p=11$ --- the
modular equation relating the multiplier to $\pprod$. What it does \emph{not} give
is the auxiliary quintic $r_{11}$ (Layer 4) or the root formula. The natural
coefficient field for $r_{11}$ is now $\Q(Y,P)$, the function field of the genus-$2$
curve $C$ of Theorem~\ref{thm:p11}, rather than $\Q(\pprod, Y)$ with $Y$ and
$\pprod$ independent --- the relation $G(Y,P) = 0$ lets us reduce powers of $Y$,
which is what makes the search for $r_{11}$ tractable. The remaining work is
therefore: find $r_{11}$ as a quintic with coefficients in $\Q(Y,P)$, decide the
exponent pattern (H1) vs.\ (H2), fix the root ordering, and evaluate at the
flagship point.

A remark on the earlier expectation is in order. Before Theorem~\ref{thm:p11} was
found, it was natural to expect the $(m,\pprod)$-relation to be a genus-$1$ plane
curve, on the grounds that $X_0(11)$ has genus $1$ and is the first prime level of
genus $1$. That expectation was wrong in two ways: the functions $m$ and $\pprod$
live on $X_0(44)$, not $X_0(11)$ (their eta-quotients need the factor $4$), and
their image curve has genus $2$, not $1$ (Remark after Theorem~\ref{thm:p11}). The
classical degree-$11$ modular equation in $\alpha,\beta$ variables (Berndt, Part
III, Entry 7) and the many degree-$11$ equations in other variables
(Vasuki--Yathirajsharma \cite{VY23}; Hart \cite{Hart08}) were never in doubt; what
was missing was specifically the equation in Ramanujan's theta variables
$(m,\pprod)$, which Theorem~\ref{thm:p11} now supplies.

\subsection{The exponent puzzle}
\label{sec:exponent}

Write $u_k = (\alpha_k^a\,\pprod/\alpha_{k+1}^b)^{1/p}$ with cyclic indices. The
proved cases give $(a,b) = (1,0)$ for $p=3,5$ (the simple form
$u_k = (\alpha_k\pprod)^{1/p}$) but $(a,b) = (2,1)$ for $p=7$ (a genuinely cyclic
form), and $(a,b) = (1,1)$ with root $1/3$ for the composite $p=9$. In all prime
cases $a-b=1$ and $\prod_k \alpha_k = \pprod^{(p+1)/2}$, which is exactly what
makes $\prod_k u_k = \pprod$. The pattern in $a$ is not determined by the proved
cases: $a = 1,1,2$ for $p=3,5,7$ does not fix $a$ for $p=11$. The two live
hypotheses are:
\begin{itemize}
\item[(H1)] \textbf{Simple form:} $u_k = (\alpha_k\pprod)^{1/11}$, i.e.\
$\alpha_k = u_k^{11}/\pprod$, and $r_{11}(\xi) = \prod_{k=1}^5(\xi-\alpha_k)$
has coefficients expressible in $(\pprod, Y)$.
\item[(H2)] \textbf{Cyclic form:} $u_k = (\alpha_k^a\pprod/\alpha_{k+1}^{a-1})^{1/11}$
for some $a \geq 2$, with $\prod \alpha_k = \pprod^6$. Then the $\alpha_k$ are
determined by cyclic sums of the $u_k$, and the coefficients of $r_{11}$ come
from those cyclic sums.
\end{itemize}
Deciding (H1) vs.\ (H2) is the first numerical experiment for the $p=11$ attack:
compute the $u_k$ to high precision at generic $q$, form $\alpha_k = u_k^{11}/\pprod$,
and test whether the elementary symmetric functions of the $\alpha_k$ lie in
$\Q(\pprod, Y)$ (H1) or whether instead the cyclic monomials $u_k^a u_{k+1}^{a-1}$
do (H2).

\subsection{Galois theory}
\label{sec:galois}

For prime $p$, $\Q(\zeta_p)^+$ has degree $(p-1)/2$ and cyclic Galois group
$(\Z/p)^\times/\{\pm 1\}$. The correct statement of the cyclicity of the auxiliary
polynomial, visible from the proved cases, is that $r_p(\xi)$ is cyclic over its
coefficient field $K = \Q(\pprod, Y)$ (Galois group $C_{(p-1)/2}$ acting by cyclic
permutation of the roots). At a special point where $\pprod, Y \in \Q$, the roots
lie in a cyclic extension of $\Q$ of degree $(p-1)/2$ --- which for $p=7$ is
$\Q(\zeta_7)^+$, but in general need not be $\Q(\zeta_p)^+$ itself. For example, at
the $p=5$ special point $n=1/5$, $\pprod = \sqrt5-1$, so the base field is
$\Q(\sqrt5) = \Q(\zeta_5)^+$, and the roots lie in the degree-$4$ field
$\Q(\sqrt5,\sqrt{10-2\sqrt5})$ --- a cyclic quadratic extension of
$\Q(\zeta_5)^+$, not in $\Q(\zeta_5)^+$ itself. For $p=11$ at $n=1/11$, the base
field is $\Q(G_{11})$ (a non-Galois cubic), so the roots of $r_{11}$ are expected to
generate a cyclic quintic extension of $\Q(G_{11})$: a degree-$15$ extension of
$\Q$, whose Galois closure is larger. This is consistent with Ramanujan's
``quintics'' label for $11$: the \emph{radical} degree needed is $5$, not that the
roots lie in $\Q(\zeta_{11})^+$.

A cyclic extension of degree $n$ is solved by Lagrange resolvents, and for
cyclotomic fields the resolvents are Gauss sums. Garrett \cite{Garrett10} treats the
quintic subfield of $\Q(\zeta_{11})$ explicitly: in $\Z[\omega]$, $\omega$ a
primitive 5th root of unity, $11 = (2+\omega)(2+\omega^2)(2+\omega^3)(2+\omega^4)$,
and the quintic subfield is generated over $\Q(\omega_5)$ by a 5th root of an
explicit product of these prime factors (Kummer--Eisenstein). Dummit's criterion
\cite{Dummit91} gives a test for a quintic to be solvable (equivalently its Galois
group contained in the Frobenius group $F_{20}$) and root formulas producing the
roots in cyclic order; this applies verbatim over the coefficient field
$K = \Q(\pprod, Y)$.

\subsection{Negative search: no low-degree relation over $\Q(\pprod,m)$}

At the flagship point $q_1 = e^{-\pi/\sqrt{11}}$ we computed
$\pprod = 0.416439434275864095678569298496$ and $m = \sqrt{11}$ exactly
(PSLQ degree $2$: $m^2 - 11 = 0$), and
\begin{equation}
u_k = 1.8352467839,\ 1.4190121186,\ 0.9294325201,\ 0.5336551405,\ 0.3223982758
\end{equation}
(strictly decreasing, as required).
With $w_k = u_k^{11}/\pprod$ we have $\prod_k w_k = \pprod^6$, as required by
either hypothesis (H1) or (H2).

We searched, by integer-relation (PSLQ) at 60 digits, for a quadratic relation
$A(\pprod,m)s_1^2 + B(\pprod,m)s_1 + C(\pprod,m) = 0$ for the power sum
$s_1 = \sum_k w_k$, with $A,B,C$ of total degree $\leq 2$ (18 monomials) and
$\leq 3$ (30 monomials). Candidate relations were found at $q=0.4$ but both
\textbf{failed verification catastrophically} at other $q$ (relative errors $6$--$322$).
They are spurious (PSLQ artefacts at 60 digits with 18--30 monomials). We conclude
that $s_1$ for $p=11$ satisfies no low-degree quadratic relation over $\Q(\pprod,m)$
with small coefficients. This is now understood: the coefficients of the auxiliary
quintic live on the genus-$2$ curve $C$, not in $\Q(\pprod,m)$ with $m$ and $\pprod$
independent --- Theorem~\ref{thm:p11} shows $Y$ and $P$ satisfy a quintic relation,
so a relation for $s_1$ over $\Q(\pprod,m)$ alone would have to absorb that
dependence. The method itself was validated by calibration on $p=7$: deriving the
$s_1$-quadratic symbolically from the known $r(\xi)$ gives a relation the
dissection's $s_1$ satisfies at $q = 0.2,\dots,0.6$ with residuals
$\leq 10^{-45}$.

\begin{remark}
The same circle of ideas explains the classical fact, noted by Bailey--Borwein--
Borwein \cite{BBB89}, that ``the first interesting prime for which an explicit
solvable form is not known is the endecadic $p=11$ case.'' An explicit solvable form
of the degree-$p$ modular equation is exactly what the theta theory needs (the
$m$--$\pprod$ equation plus the auxiliary polynomial), and $p=11$ is the first
prime where it is not known.
\end{remark}

\subsection{Attack plan for $p=11$}

We outline the concrete steps, in order of expected payoff, for completing the
$p=11$ theory. Steps 0 and 2 of the original plan (verify the product formula;
find the $(m,\pprod)$ relation) are now done --- the first by
Theorem~\ref{thm:pprod}, the second by Theorem~\ref{thm:p11}. Each remaining step
is checkable with $q$-series to 200+ digits.
\begin{enumerate}
\item \textbf{Decide (H1) vs.\ (H2)} (Step 1): test whether the elementary symmetric
functions of $\alpha_k = u_k^{11}/\pprod$ lie in $\Q(Y,P)$ (the function field of
$C$), or whether the cyclic monomials $u_k^a u_{k+1}^{a-1}$ do. Use
Theorem~\ref{thm:p11} to reduce powers of $Y$.
\item \textbf{The auxiliary quintic} (Step 2): compute the power sums of
$u_k^{11}/\pprod$ as functions on $C$ and determine $r_{11}(\xi)$; identify its
Galois group (test cyclicity $C_5$ by Dummit's sextic criterion).
\item \textbf{First test case} (Step 3): at $q = e^{-\pi/\sqrt{11}}$,
$Y = 11$ and $\pprod = 4\sqrt{2}/G_{11}^{10}$ with $G_{11}$ known in cubic radicals,
so $r_{11}$ at this point should be a quintic over $\Q(G_{11})$; compute the $u_k$
numerically, form the candidate $\alpha_k$, and determine the quintic by integer
relation, checking cyclicity by Dummit's sextic test.
\item \textbf{Root ordering} (Step 4): use $u_1 > u_2 > u_3 > u_4 > u_5$ plus a
cyclic-sum identity analogous to Reb\'ak's Lemma 3.4/3.7 to fix which cyclic ordering
of the $\alpha_k$ (5 rotations $\times$ 2 orientations $= 10$ possibilities, versus
$5! = 120$ total) gives the $u_k$ in decreasing order.
\item \textbf{Assemble the identity and the special value} (Step 5): the $p=11$ theta
identity $\varphi(q^{1/11})/\varphi(q^{11}) = 1 + \sum_{k=1}^5 u_k$ with $u_k$ as in
\S\ref{sec:exponent}, and the special value $\varphi(e^{-11\pi\sqrt{11}})$ in terms
of $\varphi(e^{-\pi\sqrt{11}})$ and 11th roots of Gaussian-period ratios.
\end{enumerate}

\section{Conclusion and open problems}
\label{sec:conclusion}

We have found and proved the degree-$11$, degree-$13$, and degree-$17$ modular
equations for Ramanujan's prime-degree theta function theory --- the first three
cases for primes $p \geq 11$ ---
together with a general product formula valid for every odd $n$, Sturm-bound proofs
of the multiplier relations of degrees $3$, $5$, and $7$, and a 60-digit numerical
value for $G_{121}$. The six proved cases suggest a uniform conjectural template
for general prime $p$ (\S\ref{sec:conjecture}), driven by the Atkin--Lehner
involution $w_4$, with testable predictions for $p = 19$ and $p = 23$; the genus
remains the one part of the template with no formula. The $p=11$ theory is not
yet complete: the auxiliary quintic
$r_{11}$, the root formula, and the flagship evaluation remain open, but
Theorem~\ref{thm:p11} supplies the modular equation that makes them tractable, and
the genus-$2$ curve $C$ is now the precise setting in which they live. The same
holds, mutatis mutandis, for $p = 13$ and $p = 17$: the auxiliary polynomials
$r_{13}$ and $r_{17}$ and their root formulas are open.

We close with the open problems that, in our view, are the most promising next steps.
\begin{openproblem}
Determine the auxiliary quintic $r_{11}(\xi)$ with coefficients in $\Q(Y,P)$,
decide between hypotheses (H1) and (H2) for the $p=11$ root formula, and fix the
root ordering.
\end{openproblem}
\begin{openproblem}
Evaluate $\varphi(e^{-11\pi\sqrt{11}})$ in terms of $\varphi(e^{-\pi\sqrt{11}})$
and 11th roots of Gaussian-period ratios --- the $p=11$ analogue of Reb\'ak's
Theorem 4.1 --- and compute $G_{1331}$ as a corollary.
\end{openproblem}
\begin{openproblem}
Compute the minimal polynomial of $G_{121}$ by CM theory (Weber modular polynomial
$\Phi_{11}$) and find a closed-form radical expression, if one exists.
\end{openproblem}
\begin{openproblem}
Find a closed form for the $w_{11}$-action on $P$ in the variables $(Y,P)$, and a
hyperelliptic model of the genus-$2$ curve $C$ of Theorem~\ref{thm:p11}.
\end{openproblem}
\begin{openproblem}
Prove Conjecture~\ref{conj:template} (the uniform template for general prime $p$,
\S\ref{sec:conjecture}) --- or determine the first prime at which it fails.
The dissection is proved, the product formula is Theorem~\ref{thm:pprod}, and the
modular equation is proved for $p = 11, 13, 17$ (Theorems~\ref{thm:p11},
\ref{thm:p13}, \ref{thm:p17}); the remaining parts are the auxiliary polynomial
of degree $(p-1)/2$ with cyclic Galois group and the cyclic $p$th-root formula.
\end{openproblem}
\begin{openproblem}
Prove the $w_4$-action \eqref{eq:w4} through the eta-multiplier (Dedekind sums)
for a representative of the Atkin--Lehner involution $w_4$ of $X_0(4p)$ --- the
single most important missing proof in the template of \S\ref{sec:conjecture},
from which the variable rule, the pure-part exponents, and the degree-$2$ map
all follow.
\end{openproblem}
\begin{openproblem}
Derive $N_{\mathrm{CM}}(p) = \#\operatorname{Fix}(w_4) - 2$ from the CM
fixed-point theory of $w_4$ on $X_0(4p)$, turning the genus in
Conjecture~\ref{conj:template} from a minimal-$R'$ guess into a prediction; in
particular, predict $g(C)$ for $p = 19$ and $p = 23$ before the Sturm
computations are run.
\end{openproblem}
\begin{openproblem}
Extend the Sturm-bound prover to $\Gamma(2)$ (half-integer eta products and
resultant elimination) to complete the proof of Schr\"oter's degree-$11$ modular
equation. (The theta-variable degree-$17$ equation is now
Theorem~\ref{thm:p17}; the classical degree-$17$ equations of Entry 12 in
$(\alpha,\beta)$ variables are a separate matter.)
\end{openproblem}

\section*{Acknowledgments}
The author thanks Bruce C.~Berndt, \H{o}rs Reb\'ak, and Sun Kim for the papers on
which this work builds, and the Ramanujan community for keeping the notebooks alive.
The numerical computations were carried out in Python with \texttt{sympy} and
\texttt{mpmath}; the author used AI assistance for code and drafting, with all
mathematical claims verified independently against the cited sources and by
recomputation. All scripts are available from the author on request.

\begin{thebibliography}{99}

\bibitem{BR25} B.~C.~Berndt and \H{o}.~Reb\'ak, \emph{Cubic and quintic analogues of
Ramanujan's septic theta function identity}, Ramanujan J. \textbf{66} (2025), art.~66.
DOI: 10.1007/s11139-024-01013-2. arXiv:2312.15513.

\bibitem{Reb23} \H{o}.~Reb\'ak, \emph{The three missing terms in Ramanujan's septic
theta function identity}, Ramanujan J. \textbf{60} (2023), 885--911.
DOI: 10.1007/s11139-022-00584-2. arXiv:2112.11884.

\bibitem{KR25} S.~Kim and \H{o}.~Reb\'ak, \emph{Ninth degree analogue of Ramanujan's
septic theta function identity}, arXiv:2506.01181 (2025).

\bibitem{Son01} S.~H.~Son, \emph{Septic theta function identities in Ramanujan's lost
notebook}, Acta Arith. \textbf{98} (2001), 361--374.

\bibitem{BerndtIII} B.~C.~Berndt, \emph{Ramanujan's Notebooks, Part III}, Springer,
1991. DOI: 10.1007/978-1-4612-0965-2.

\bibitem{BBB89} D.~H.~Bailey, J.~M.~Borwein, and P.~B.~Borwein, \emph{Ramanujan,
modular equations, and approximations to $\pi$}, Amer. Math. Monthly \textbf{96}
(1989), 201--219.

\bibitem{Hart08} W.~B.~Hart, \emph{Schl\"afli modular equations for generalized Weber
functions}, Ramanujan J. \textbf{15} (2008), 435--468.
DOI: 10.1007/s11139-007-9087-8.

\bibitem{Garrett10} P.~Garrett, \emph{Kummer theory and Eisenstein series}, 2010.
\url{https://www-users.cse.umn.edu/~garrett/m/v/kummer_eis.pdf}

\bibitem{Dummit91} D.~S.~Dummit, \emph{Solving solvable quintics}, Math. Comp.
\textbf{57} (1991), 387--401.

\bibitem{VY23} K.~R.~Vasuki and M.~V.~Yathirajsharma, \emph{Crossed out modular
equations of degree 11 of Ramanujan}, Ramanujan J. \textbf{61} (2023), 1145--1150.
DOI: 10.1007/s11139-022-00630-z.

\bibitem{BR21} B.~C.~Berndt and \H{o}.~Reb\'ak, \emph{Explicit values for Ramanujan's
theta function $\varphi(q)$}, Hardy--Ramanujan J. \textbf{44} (2021), 41--50.
DOI: 10.46298/hrj.2022.8923. arXiv:2112.11882.

\end{thebibliography}

\end{document}
